% Commentary on manuscript pages 1--18.

\begin{dcomment}{exposition}{p1}{A combinatorial interval carries two orders}
Grothendieck begins again because a path should be easy to concatenate and easy to change without losing its endpoints. The extra conceptual effort announced here consists in making its shape an object of study. Before forming a category of all paths, he isolates the finite categories which can serve as their domains. The commentary on Chapter~VI, manuscript p.~69 (``Concatenation and the two iterated units'') discusses the preceding path construction; the present chapter reorganizes its combinatorics.

A standard interval is the free category on an oriented path graph. Write its transition signs as
\[
 \tau=(\varepsilon _1,\ldots,\varepsilon _n),\qquad
 \varepsilon_i\in\{+,-\}.
\]
There is one object at each integer from $0$ to $n$. A positive sign supplies the generating arrow $i-1\to i$ and a negative sign supplies $i\to i-1$. The free category is thin: between two fixed vertices there is at most one directed path, since the underlying undirected graph has no cycle. Thus the description as a category associated with a partially ordered set introduces no additional relations. Conversely, the two conditions in the text reconstruct precisely this free category.

For later use we record the complete hom-set calculation. If $i<j$, then
\[
\begin{split}
 I_\tau(i,j)\ne\varnothing
   &\quad\Longleftrightarrow\quad
   \varepsilon_{i+1}=\cdots=\varepsilon_j=+,\\
 I_\tau(j,i)\ne\varnothing
   &\quad\Longleftrightarrow\quad
   \varepsilon_{i+1}=\cdots=\varepsilon_j=-.
\end{split}
\]
Every nonempty hom-set contains exactly one arrow. These formulas prove both the classification by sign words and the count $2^n$. They also explain why a functor out of an interval will be specified freely by its values on the transition arrows.

The numerical order and the categorical partial order have different jobs. The first tells us in which order the vertices are visited; the second tells us which composites exist. For the shape $0\to1\leftarrow2$, the designated source of the interval is $0$ and its designated target is $2$, although neither statement says that $0$ is initial or $2$ terminal. A path of this shape travels along the first arrow and against the second. Confusing these two orders would already make the later distinction between an opposite path and an opposite category disappear.

Length counts transitions. The interval of length zero has one vertex and no transition, and the interval of length $n$ has $n+1$ vertices. We keep this convention throughout; it will identify the isolated equality $n=\operatorname{card}I$ on ms.~13 as a printed slip rather than a competing convention.
\end{dcomment}

\begin{dcomment}{exposition}{p4}{Runs, simplices, and the meaning of maximality}
Grothendieck next separates the places where a path can be composed from the places where its direction turns. A maximal constant-sign run of the word $\tau$ is his structural chain. If the run lengths are $l_1,\ldots,l_r$, then
\[
 n=l_1+\cdots+l_r,
\]
and the sign of the $a$th run is $(-1)^{a-1}\varepsilon_1$, when signs are written as $\pm1$. This proves the parity rule for initial and final directions. The useful length $r$ records the number of changes of compositional regime, rather than the number of arrows which happen to have been written down.

The paragraph identifying chains with simplex subcategories needs a small but substantive separation. A chain was defined on ms.~3 to be a contiguous set of transitions. An object-injective functor $\Delta_k\to I$ may skip vertices. For example,
\[
 \Delta_1\longrightarrow\Delta_2,
 \qquad 0\longmapsto0,\quad1\longmapsto2,
\]
is fully faithful, but its image is not the domain of one transition of $\Delta_2$; nor does it describe its maximal structural chain, whose length is two. The words ``chaînes maximales'' and the unrestricted injectivity condition both occur in the edition. They therefore remain in the transcription, with their correction supplied here.

The exact statement is useful in its own right. For $k>0$, an injective functor from $\Delta_k$ picks $k+1$ vertices in a single structural interval, in their categorical order. Indeed its first and last image vertices are connected by an arrow, so every vertex between them in the numerical order belongs to the same constant-sign run. All intermediate image vertices lie between these two in that run's categorical order. Conversely, any such selection determines an injective functor. It is fully faithful because the target is thin and its order on those vertices is exactly the simplex order. Its object function is strictly increasing in the external orientation for a positive run and strictly decreasing for a negative run.

To recover a chain in the earlier sense, impose numerical contiguity of the image. To recover a maximal structural chain, require the image to be the whole structural interval. These distinctions do not disturb the intended decomposition into runs. They identify which extra condition the printed correspondence omitted, while preserving its correct full-faithfulness assertion.
\end{dcomment}

\begin{dcomment}{exposition}{p5:Le groupe des automorphismes}{The endpoint determines an orientation, not an endofunctor}
The purpose of Proposition~1 is to make orientation intrinsic enough to transport it through isomorphisms. Its first assertion is correct. At an endpoint, one of the two slices is a point and the other is the linearly ordered category running along the unique adjacent structural chain. An interior vertex at which the signs agree has both a nonidentity incoming and a nonidentity outgoing arrow, so neither slice is a point. At a turning vertex, one slice is a point, but the other has two nonempty branches meeting at their common source or common target. In the source case it has two incomparable maximal objects, and in the target case two incomparable minimal objects. It cannot be a simplex category. This gives the stated characterization of the two endpoints.

The second assertion is false with the word ``foncteur'' which the edition prints. On $\Delta_2$ the monotone map
\[
 0\longmapsto0,\qquad1\longmapsto0,\qquad2\longmapsto2
\]
defines a nonidentity endofunctor fixing both endpoints. The arrow $1\to2$ is sent to the composite $0\to2$; there is no requirement that a functor preserve a generating transition as a generating transition. This is exactly the flexibility which the later category of types is intended to accommodate.

What the subsequent argument proves is that an \emph{automorphism} fixing the endpoints is the identity. Here is a direct proof which also supplies the uniqueness statements used on ms.~7--9. A nonidentity arrow of an interval is indecomposable into two nonidentity arrows if and only if it is a transition: the displayed hom-set formula proves this, by factoring every longer directed segment at an intermediate vertex. Isomorphisms preserve indecomposability. The graph of indecomposable arrows, after forgetting directions, is the ordinary path with vertices $0,\ldots,n$. A graph automorphism fixing $0$ fixes its unique neighbour, then that neighbour's still unvisited neighbour, and so on. It fixes every vertex and hence every arrow. Every graph automorphism is consequently the identity or vertex reversal. The categorical automorphism group is a subgroup of this two-element group.

This is also the intrinsic reconstruction of the external total order: choose one endpoint, call it $s_0$, and successively take the unique neighbour not already visited. Grothendieck's slice argument implements the same reconstruction in categorical terms. In its inductive paragraph the slices must be taken in the remaining full tail $I'$ at $s_i$; the printed $i\backslash I$ and $I/i$ have dropped that qualification. Neither reconstruction is preserved by arbitrary functors, which may collapse or lengthen transitions.

The unfinished sentence about characterizing the \emph{source} of the interval should not be completed by guessing. The underlying category determines the pair of endpoints, but a symmetric interval has an automorphism interchanging them. Selecting one of them is precisely the additional orientation. After the definition of a reduced interval on ms.~12 we prove a sharp version of the stronger rigidity assertion: all endpoint-fixing endofunctors are identities exactly for reduced intervals.
\end{dcomment}

\begin{dcomment}{examples}{p9}{Reversal and categorical duality in a sign word}
Grothendieck's two uses of ``opposite'' describe operations which can be written down independently. This makes the examples of symmetry and self-duality entirely explicit. For a sign word $\tau=(\varepsilon_1,\ldots,\varepsilon_n)$, put
\[
\begin{aligned}
 R\tau&=(-\varepsilon_n,\ldots,-\varepsilon_1),\\
 D\tau&=(-\varepsilon_1,\ldots,-\varepsilon_n).
\end{aligned}
\]
The first is the standard model of the same category with its external orientation reversed; the standardizing isomorphism sends $i$ to $n-i$. The second is the categorical opposite with the original numerical orientation. The two involutions commute, and $RD\tau$ is the word read backwards without changing its signs.

For $n>0$, the interval is symmetric precisely when
\[
 \varepsilon_i=-\varepsilon_{n+1-i}\qquad(1\leq i\leq n).
\]
Necessity follows because a nonidentity automorphism must be reversal, by the preceding proof; sufficiency follows by checking each transition. This condition forces $n$ to be even. Likewise $I\cong I^o$ precisely when
\[
 \varepsilon_i=\varepsilon_{n+1-i}\qquad(1\leq i\leq n).
\]
The isomorphism must reverse the numerical order: an order-preserving bijection of the vertex line would be the identity and would have to reverse a nonidentity transition without changing its endpoints. These two sign conditions cannot hold together for a nonempty word.

For instance, $(+,-)$ and $(-,+)$ satisfy the first condition. They give the cospan and span, each symmetric, but the two categories are not isomorphic: one has a terminal object and two minimal objects, the other an initial object and two maximal objects. The constant word $(+,+,+)$ satisfies the second condition and gives the self-dual category $\Delta_3$ with no nonidentity automorphism. A less uniform self-dual example is $(+,-,+)$; its middle sink and source are exchanged under the anti-isomorphism. For a reduced alternating word, symmetry holds at positive even length and self-duality at odd length.

The point is separate: it has a unique orientation and only the identity automorphism, and both operations leave it unchanged. Keeping this case explicit avoids construing its coincident endpoints as two different choices.
\end{dcomment}

\begin{dcomment}{exposition}{p10}{Critical vertices and the correct deletion test}
Grothendieck wants to retain exactly the vertices which a path cannot bypass by composing arrows. These are the endpoints of the structural intervals. If their positions are
\[
 0=c_0<c_1<\cdots<c_r=n,
\]
then $[c_{a-1},c_a]$ is one constant-sign run, and each interior $c_a$ is a change of sign. In the categorical partial order, the critical vertices are exactly the minimal and maximal vertices. This observation explains the two slices in Proposition~3, but it also exposes an error in its global formulation.

Let
\[
 A=(I/s)\setminus\{s\},\qquad
 B=(s\backslash I)\setminus\{s\}.
\]
These are respectively the strict predecessors and strict successors of $s$ in the categorical order. They do not in general exhaust the other vertices. In $0\to1\leftarrow2$, take $s=0$. It is critical; $A$ is empty and $B$ consists of $1$, whereas $I\setminus\{0\}$ still contains $2$. Thus the printed equality with all of $I\setminus\{s\}$ cannot be right. The edition prints the same reading.

The footnote gives the valid local test: there is no arrow between $A$ and $B$. If $s$ is critical, it is minimal or maximal and one of these categories is empty. If $s$ is not critical, its two adjacent transitions have the same sign. One adjacent vertex is a strict predecessor and the other a strict successor; their composite through $s$ is an arrow between $A$ and $B$. Equivalently, the full subcategory on the vertices comparable with $s$, after removing $s$, is the coproduct $A\amalg B$ exactly when $s$ is critical.

For the application to subintervals on the next page, the appropriate \emph{global} test uses the numerical order. Suppose $0<s<n$, and let $L$ and $H$ be the full subcategories on $[0,s-1]$ and $[s+1,n]$. Then
\[
 I\setminus\{s\}=L\amalg H
 \quad\Longleftrightarrow\quad s\text{ is critical}.
\]
If $s$ is a turning vertex, an arrow from one side to the other would traverse two transitions of opposite signs, which is impossible. If it is not a turning vertex, the two neighbouring vertices are connected by their composite through $s$, an arrow still present in the full subcategory obtained by deleting $s$. This corrected statement proves the separation property used in Proposition~5. At a numerical endpoint, deleting it leaves a connected interval; its criticality is already accounted for by its being an endpoint.
\end{dcomment}

\begin{dcomment}{examples}{p11}{Connected full subcategories and the vertices they must retain}
Grothendieck is preparing a reduction which composes consecutive arrows without identifying a change of direction. The distinction between a full subcategory and an arbitrary subcategory matters here. A full subcategory retains every composite whose endpoints remain, even when an intermediate vertex has been removed. Thus numerical gaps do not necessarily disconnect it.

Let $J=\{j_0<\cdots<j_m\}$ be a nonempty subset of the vertices of $I$, with its full induced category. Two successive chosen vertices $j_a,j_{a+1}$ are joined by an arrow if and only if the signs between them are constant. By Proposition~4 this is equivalent to there being no critical vertex strictly between them. Whenever this condition fails, there can be no arrow from any chosen vertex at or before $j_a$ to any chosen vertex at or after $j_{a+1}$, in either direction: such an arrow would cross the same change of sign. The connected components are therefore the maximal consecutive blocks of the list $(j_0,\ldots,j_m)$ in which successive entries are directly connected.

Within one such block, make a generating transition between each successive pair, in the direction of its arrow in $I$. The free category on this oriented line is exactly the induced category. To verify fullness, an arrow between nonconsecutive retained vertices exists in $I$ precisely when the signs on the whole intervening segment agree. That is also precisely the condition that the successive arrows of the block compose. This proves Proposition~5(a), including the claim about the induced orientation. It proves (b) as well: the whole chosen list is connected exactly when every ambient critical vertex between its extreme vertices was retained.

For a concrete calculation take
\[
 I:\quad0\longrightarrow1\longrightarrow2
       \longleftarrow3\longleftarrow4\longrightarrow5.
\]
Its critical vertices are $0,2,4,5$. The full category on $\{0,2,4,5\}$ is the reduced interval
\[
 0\longrightarrow2\longleftarrow4\longrightarrow5.
\]
Deleting $1$ or $3$ has merely replaced a pair of arrows by its composite. The full category on $\{0,1,4,5\}$ has two components, $\{0,1\}$ and $\{4,5\}$, because the omitted turning vertex $2$ separates them. The least subinterval containing $1$ and $5$ is $\{1,2,4,5\}$: the endpoints must remain, the two intervening turns must remain, and the preceding calculation shows that nothing else is required. This is the corollary on ms.~12.

The last sentence of footnote~4 asserts more than fullness permits: an interval subcategory need not be full. Inside $\Delta_2$, keep all three objects and the arrows $0\to1$ and $0\to2$, but omit $1\to2$. This is a subcategory, because its two nonidentity arrows are not composable. Abstractly it is the span interval with vertex order $1,0,2$. It is not full in $\Delta_2$. Thus fullness must remain part of the definition of a subinterval; it cannot be inferred solely from the abstract shape of the subcategory.

Finally, for a full subinterval with at least two vertices, an interior retained vertex is a turn in the induced shape exactly when it was an ambient critical vertex. This proves Proposition~6. The reduced interval is consequently the smallest full connected part of $I$ which contains its endpoints, and its length is the useful length of $I$.
\end{dcomment}

\begin{dcomment}{development}{p12}{The critical-vertex inclusion has two adjoints}
The reduced interval retains every turn, but Grothendieck has not yet described how to return to it from an arbitrary vertex. There are two particularly useful choices, one following the categorical source of a run and one following its categorical target. They relate reduction to the endpoint-preserving transformations which a path calculus needs.

\textbf{Two canonical collapses.}\ours{}
Let $j:R=I_{\mathrm{red}}\hookrightarrow I$ be the full inclusion of the critical vertices of an oriented interval. There are retractions $r_s,r_t:I\to R$ with
\[
 r_t\dashv j\dashv r_s,\qquad
 r_sj=r_tj=\mathrm{id}_R.
\]
Both retractions fix the endpoints. At a noncritical vertex $x$, $r_s(x)$ is the categorical source of its structural interval and $r_t(x)$ is that interval's categorical target. At a critical vertex both are the identity. For the point we set $R=I$ and take all three functors to be the identity.

\emph{Proof.} A noncritical vertex lies strictly inside a unique structural interval. On a positive run $[a,b]$, the source-collapse sends every interior vertex to $a$, leaves $b$ at $b$, and is a monotone map of its categorical chain. On a negative run the source is $b$ and the same argument applies with the categorical order reversed. The target-collapse is obtained by exchanging source and target. The run-wise maps agree at shared critical vertices. Every arrow belongs to one run, so these formulas give functors on the whole interval.

To verify the adjunctions rather than merely exhibit homotopies, let $q$ be a critical vertex. If $x$ is itself critical, fullness of $j$ gives the required hom-set bijections. If $x$ is interior to a run, the only critical vertices comparable with $x$ are its source $a$ and target $b$, with $a\to x\to b$ in the categorical order. Indeed an arrow extending beyond either endpoint of that run would have to pass a change of sign. Consequently
\[
\begin{split}
 I(jq,x)\ne\varnothing&\ \Longleftrightarrow\ q=a
     \ \Longleftrightarrow\ R(q,r_sx)\ne\varnothing,\\
 I(x,jq)\ne\varnothing&\ \Longleftrightarrow\ q=b
     \ \Longleftrightarrow\ R(r_tx,q)\ne\varnothing.
\end{split}
\]
For the middle equivalences we use that a source endpoint is a minimal vertex of $I$ and a target endpoint a maximal vertex; hence the same holds in $R$. All these hom-sets are empty or singletons, so the bijections are unique and natural. They establish $j\dashv r_s$ and $r_t\dashv j$. The two retraction identities follow directly from the definition. The counit and unit are the unique arrows
\[
 j r_s(x)\longrightarrow x,\qquad x\longrightarrow j r_t(x),
\]
which are identities whenever $x$ is critical. This also verifies the triangle identities directly.\hfill$\square$

These natural transformations are categorical deformations relative to all critical vertices, in particular relative to both endpoints. Applying the nerve, a natural transformation $F\Rightarrow G$ supplies a simplicial homotopy through the functor $I\times\Delta_1\to I$ which it defines. Thus $Nj$ has homotopy inverses $Nr_s$ and $Nr_t$. Nothing in this argument asserts that the interval categories themselves are equivalent: a skeletal interval with an interior noncritical vertex has strictly more objects than its reduction.

The effect on diagrams gives a more informative comparison. For every category $X$, precomposition yields
\[
 r_s^*\dashv j^*\dashv r_t^*,\qquad
 j^*r_s^*=j^*r_t^*=\mathrm{id}_{\operatorname{Fun}(R,X)}.
\]
This reversal of adjunctions under precomposition is the standard calculation of \citet[Proposition~4.3.7]{Riehl2016}; it can also be read directly from the two displayed natural transformations by applying each diagram $I\to X$. The source extension inserts identity arrows before the final arrow of each directed run, and the target extension inserts them after its first arrow. Restriction to $R$ composes each run.

For $I=\Delta_2$, $R=\{0,2\}$ and $X=\Delta_1$, consider the two diagrams with vertex values
\[
 F=(0,0,1),\qquad G=(0,1,1).
\]
They have the same restriction $(0,1)$ to $R$. There is no natural transformation $G\to F$, since its middle component would have to be $1\to0$ in $\Delta_1$, whereas their restrictions have an identity map. Hence $j^*$ is not full and is not an equivalence. If $X$ is a groupoid, every component of both comparison transformations is invertible, and $j^*$ is an equivalence, with either section as an inverse. Reduction therefore loses factorization data in general, but loses no such data when all arrows of the target are already invertible.

The general precomposition-adjunction statement is known and cited above. The critical-vertex reduction itself also appears in \citet[Lemma~4.16]{ArakawaCnossen2026}. For the additional explicit adjoints and the sharp rigidity criterion below, we checked that construction, the local interval arguments here, and the adjunction calculus in \citet[\S4.3]{Riehl2016}; targeted searches for critical-vertex reductions and their adjoints did not locate these particular formulations. We supply the extra deductions as elementary developments of the present definitions, without a priority claim.
\end{dcomment}

\begin{dcomment}{development}{p12}{The exact scope of endpoint rigidity}
The false generality of Proposition~1(b) suggests asking precisely which part of an interval allows an endpoint-preserving functor to move a vertex. The preceding collapses give one implication, while a finite displacement argument gives the other.

\textbf{Endpoint rigidity criterion.}\ours{}
For a finite oriented interval $I$, the following are equivalent: $I$ is reduced; every endofunctor of $I$ fixing its two designated endpoints is the identity. We include the point as a reduced interval.

\emph{Proof.} Suppose $I$ is reduced, of length $n>0$, and identify its vertices with $[0,n]$. A nonidentity arrow of $I$ is then a single transition: two consecutive transitions have opposite signs and cannot compose. If $f:I\to I$ is any functor, the image of a transition is therefore either an identity or a transition. Its object function satisfies
\[
 -1\leq f(i)-f(i-1)\leq1\qquad(1\leq i\leq n).
\]
If $f(0)=0$ and $f(n)=n$, summation gives
\[
 n=f(n)-f(0)=\sum_{i=1}^n\bigl(f(i)-f(i-1)\bigr)\leq n.
\]
Equality forces every summand to be $1$. Thus $f(i)=i$ for all $i$, and thinness forces $f$ to be the identity on arrows as well. We have made no monotonicity assumption on $f$; it is a consequence of the equality. The point case is immediate.

Conversely, if $I$ is not reduced, a structural interval has an interior vertex. The endofunctor $j r_s$ of the preceding result fixes every critical vertex and hence both endpoints, but sends that interior vertex to the source of its run. It is not the identity.\hfill$\square$

The obstruction is therefore exactly a nontrivial factorization inside a constant-sign run. Automorphisms cannot change such a factorization because they preserve the graph of indecomposable arrows. General functors can replace it by an identity followed by a longer composite. This is why the automorphism proof on ms.~6--7 is sound while its printed extension to all endofunctors fails.
\end{dcomment}

\begin{dcomment}{exposition}{p14}{A fixed path type is a freely generated diagram}
Grothendieck now uses the work on intervals to make path data and path morphisms completely explicit. Fixing the type does not impose any compatibility among the arrows of the path beyond their sources and targets. This is the benefit of having chosen a free path category as the domain.

For example, let the type be $(+,+,-)$. A diagram consists of
\[
 a_0\xrightarrow{u_1}a_1\xrightarrow{u_2}a_2
       \xleftarrow{u_3}a_3.
\]
The only nonidentity arrow of the domain which is not a transition is $0\to2$, and its image must be $u_2u_1$. Given a second diagram $a'_0\to a'_1\to a'_2\leftarrow a'_3$, a map between them consists of four arrows $\alpha_i:a_i\to a'_i$ satisfying
\[
 \alpha_1u_1=u'_1\alpha_0,\qquad
 \alpha_2u_2=u'_2\alpha_1,\qquad
 \alpha_2u_3=u'_3\alpha_3.
\]
Naturality for the longer arrow follows by calculation:
\[
 \alpha_2(u_2u_1)
 =(u'_2\alpha_1)u_1
 =u'_2(u'_1\alpha_0)
 =(u'_2u'_1)\alpha_0.
\]
Every longer constant-sign run is handled by the same induction, with the index order reversed when its arrows are retrograde. Thus the displayed transition squares are both necessary and sufficient for a natural transformation. Identities and composition are taken vertex by vertex; composing commutative squares verifies that the result is again a morphism of paths.

This also gives a useful exact model of the fixed-type category. Let $\operatorname{Arr}(X)=\operatorname{Fun}(\Delta_1,X)$, with ordinary domain and codomain functors $d,c$. For a positive transition its ordered endpoint pair is $(d,c)$; for a negative one it is $(c,d)$. Then $\operatorname{Fun}(I_\tau,X)$ is the iterated strict fibre product of $n$ copies of $\operatorname{Arr}(X)$, matching the second ordered endpoint of one to the first ordered endpoint of the next. On objects this identifies the shared vertices. On arrows it identifies the corresponding vertical components of the naturality squares. The two descriptions are therefore isomorphic as categories, not merely bijective on objects. At length zero the category is $X$ itself.

The construction retains all factorizations and all commuting ladders between them. It has not inverted the retrograde arrows in $X$ and has not identified paths related by a ladder. Those operations would require additional choices or a localization. Here a retrograde transition simply says that the arrow of $X$ points against the chosen order of traversal. Keeping this distinction explicit will matter when the category of types starts to vary.
\end{dcomment}

\begin{dcomment}{context}{p16:sont opposées l'une de l'autre}{Zigzags before taking a localization}
Grothendieck's paths belong to the same categorical language as fractions, while retaining information which a category of fractions has already identified. In \citet[Chapter~I, \S1.1]{GabrielZisman1967}, localization adjoins inverses to a specified collection of arrows and imposes the relations needed for functoriality and inversion. A span $x\leftarrow z\to y$ can consequently represent a morphism from $x$ to $y$ when its left arrow is among those to be inverted. The fixed-type path category written here makes no such selection: the objects and every component of a natural transformation are ordinary categorical data.

The two-dimensional comparison is equally concrete. The reduced hammocks of \citet[Definition~2.1]{DwyerKan1980Calculating} consist of commutative ladders of zigzags, but their vertical arrows and their left-pointing arrows are required to lie in a specified subcategory of weak equivalences; identity columns are removed and neighbouring columns of the same direction are composed. Their Proposition~3.1 identifies connected components of the resulting mapping spaces with ordinary localized morphisms. Grothendieck's full functor category of one fixed interval is therefore an ingredient of a comparable calculus, not by itself a hammock mapping space. The later commentary on the varying category of types explains the precise relation; here the historical comparison tells us why it is useful to retain ladders rather than only strings of arrows.
\end{dcomment}

\begin{dcomment}{examples}{p16}{A span and a cospan give different diagram categories}
Grothendieck emphasizes that the direction of a turn changes the fixed-type category. The difference can be seen without any localization or homological algebra. We give a finite example which also distinguishes equivalence of categories, not only strict isomorphism.

Let $X$ be the three-object poset
\[
             a\longleftarrow b\longrightarrow c.
\]
A span diagram in $X$ is a triple $(x_0,x_1,x_2)$ with $x_1\leq_X x_0$ and $x_1\leq_X x_2$. If its middle value is $b$, each endpoint can independently be $a,b,$ or $c$, giving nine diagrams. If the middle value is $a$, both endpoints must be $a$; if it is $c$, both must be $c$. Thus
\[
 \bigl|\operatorname{Ob}\operatorname{Fun}(\Phi_1,X)\bigr|
 =3^2+1^2+1^2=11.
\]
A cospan diagram has $x_0\leq_X x_1$ and $x_2\leq_X x_1$. For middle value $b$ there is one possibility. For middle value $a$, both endpoints independently lie in $\{b,a\}$, giving four; for middle value $c$, both independently lie in $\{b,c\}$, again giving four. Hence
\[
 \bigl|\operatorname{Ob}\operatorname{Fun}(\Psi_1,X)\bigr|
 =1^2+2^2+2^2=9.
\]
Morphisms in both functor categories are completely determined: there is one from $(x_0,x_1,x_2)$ to $(y_0,y_1,y_2)$ exactly when $x_i\leq_Xy_i$ for each $i$. The naturality squares then commute automatically because $X$ is thin. These categories are posets and have no distinct isomorphic objects. An equivalence between posets is an isomorphism, since essential surjectivity is actual surjectivity and full faithfulness reflects their order. Their different object counts therefore rule out an equivalence.

Orientation reversal is a separate operation. For the span, interchanging $x_0$ and $x_2$ gives another span and an automorphism of its diagram category. It leaves the middle vertex and both arrow directions of the underlying span shape intact, while exchanging the two endpoint evaluations. Passing from the span to the cospan reverses the categorical arrows at the turn. The preceding count detects exactly that change.
\end{dcomment}

\begin{dcomment}{later}{p16}{When a change of orientation becomes a stable equivalence}
The span--cospan distinction remains real for the ordinary functor categories which Grothendieck has just defined. A later theorem explains why analogous diagram categories can nevertheless become equivalent after one passes to a stable homotopy theory. For a finite oriented tree $Q$, \citet[Theorem~9.19 and Corollary~9.23, arXiv v3]{GrothStovicek2016Trees} construct reflection equivalences and deduce strong stable equivalence between any two orientations of the same underlying tree. In particular, for a stable derivator $\mathbb D$,
\[
 \mathbb D^{\Phi_1}\simeq\mathbb D^{\Psi_1}.
\]
These are equivalences of shifted derivators, natural with respect to exact morphisms. They are constructed using stable operations such as fibres, cofibres, and biproducts; they are not induced by an isomorphism of the span and cospan categories. The finite example above therefore presents no obstruction. It concerns arbitrary ordinary targets, whereas the reflection theorem uses stability to transform the values and arrows of the diagram. This gives a precise later setting in which changing the signs of an interval preserves a diagram theory even though it does not preserve its indexing category.
\end{dcomment}

\begin{dcomment}{exposition}{p17}{The endpoints are part of the fixed-type object}
Grothendieck's insistence on orientation is most transparent over the category of endpoint pairs. The category attached to an oriented type $(I,\omega)$ is not merely $\operatorname{Fun}(I,X)$; it comes with the functor
\[
 (\sigma,\beta):\operatorname{Fun}(I,X)\longrightarrow X\times X.
\]
On a natural transformation $\alpha:c\Rightarrow d$, its value is $(\alpha_0,\alpha_n)$. Thus a morphism of paths may move the endpoints, and the displayed functor records exactly how. Its strict fibre over $(x,y)$ consists of paths with endpoints $x,y$ and natural transformations which are the identity at both endpoints. This fibre is a concrete category of fixed-endpoint ladders.

Let $\omega'$ be the reverse orientation on the same underlying interval category. The identity on that underlying category identifies the two ordinary functor categories; under this identification the endpoint pair changes from $(c(0),c(n))$ to $(c(n),c(0))$. Consequently the comparison is a functor over the transposition $s(x,y)=(y,x)$. If $x\ne y$, it carries the fibre over $(x,y)$ to the fibre over $(y,x)$, not to itself. This explains why equality of the underlying functor categories does not erase the distinction between a path and its reverse.

The endpoint map $\alpha:\{0,1\}_{\mathrm{disc}}\to I$ packages the same structure as restriction. For positive length its two images are distinct. For length zero it is the fold onto the sole vertex, so the word ``inclusion'' on the next page needs that qualification; restriction is then the diagonal $X\to X\times X$. All the formulas still hold. The subsequent discussion of symmetric types compares this identity-on-underlying-diagram reversal with the additional comparison induced by a reflection of the shape. Their difference belongs to the structure over $X\times X$, which is why the base category has been kept visible.
\end{dcomment}

\begin{dcomment}{exposition}{p20}{Operations on paths are maps of their shapes}
Grothendieck wants to do the combinatorics once, in a category small enough that every operation can be seen. A path in an arbitrary category carries many labels, but an operation valid for every category cannot depend on accidental properties of those labels. Its entire content is already present in the universal path. The passage from an oriented interval to its sign word makes this universal object unambiguous. It does not assert that a particular category of paths has no automorphisms; it removes the ambiguity coming from presenting the same unoriented interval in two different ways.

Write $P_I(X)=\Hom_{\Cat}(I,X)$ for the \emph{set} of paths of shape $I$, keeping $X^I=\uHom(I,X)$ for their category. The distinction between these two expressions matters in the displayed Yoneda argument. If
\[
 L_X:P_I(X)\longrightarrow P_J(X)
\]
is natural in all functors $X\to Y$, put $\varphi=L_I(\mathrm{id}_I):J\to I$. Applying naturality to the functor $c:I\to X$ gives
\[
 L_X(c)=c\circ L_I(\mathrm{id}_I)=c\varphi.
\]
Thus every such operation is precomposition with exactly one $\varphi$, and every $\varphi$ defines such an operation. This is the covariant Yoneda lemma, with the representing objects appearing in the reverse order; see \citet[Theorem~2.2.4 and Corollary~2.2.8]{Riehl2016}. No hypothesis about interval categories is needed for this argument.

There is then a canonical operation on morphisms of paths. If $v:c\Rightarrow d$, set $(L_Xv)_j=v_{\varphi(j)}$. For an arrow $a:j\to k$ of $J$, its naturality square is the naturality square of $v$ at $\varphi(a)$. Identities and composites are preserved component by component. Consequently the operation lifts to the functor
\[
 \varphi^*:X^I\longrightarrow X^J.
\]
This proves the extension asserted on manuscript page~23, without assuming a separate coherence theorem. Natural transformations between changes of shape have an equally concrete meaning: $\alpha:\varphi\Rightarrow\psi$ gives $c\alpha:c\varphi\Rightarrow c\psi$, naturally in $c$ and $X$.

The special advantage of intervals now enters. They are posets, so a functor into an interval is determined by its object function, and a natural transformation between two such functors exists precisely when the corresponding pointwise arrows exist. If it exists, it is unique. This uniqueness makes every required diagram of \emph{shape transformations} commute. It does not make the categories $X^I$ thin when $X$ is an arbitrary category: those categories still remember all morphisms of $X$.

For fixed endpoints the same universal-path argument is even more economical. Suppose $I$ and $J$ have specified sources and targets. The operation $c\mapsto c\varphi$ preserves both endpoints for every $X$ and $c$ exactly when $\varphi(0_J)=0_I$ and $\varphi(1_J)=1_I$. Necessity follows by taking $X=I$ and $c=\mathrm{id}_I$; sufficiency is direct evaluation. This is the reason the first category of all changes of shape will soon have to be restricted.
\end{dcomment}

\begin{dcomment}{examples}{p23:entre les ensembles d'objets correspondants.}{Composition and insertion of an identity}
The two examples on manuscript pages~21--22 illustrate how a familiar operation becomes a finite calculation before any coefficient category is chosen. Take $I=[2]=(0\to1\to2)$ and $J=[1]$. The unique endpoint-preserving functor $\varphi:J\to I$ sends $0$ to $0$ and $1$ to $2$. A path $c=(a_0\xrightarrow{u}a_1\xrightarrow{v}a_2)$ therefore becomes the arrow $vu:a_0\to a_2$.

A morphism from $c$ to $d=(b_0\xrightarrow{u'}b_1\xrightarrow{v'}b_2)$ is a triple $(h_0,h_1,h_2)$ satisfying $h_1u=u'h_0$ and $h_2v=v'h_1$. The induced morphism between the composite arrows is the square with sides $h_0,h_2$. Its commutativity follows from the explicit computation
\[
 h_2vu=v'h_1u=v'u'h_0.
\]
The middle component has done its work and can be forgotten. The reversed variance is essential: an operation which reduces the length of a path is represented here by a functor from the shorter interval to the longer one.

For the second operation let $J=I=[2]$ and take the nondecreasing object function $\psi=(0,2,2)$. Then $\psi^*c=(a_0\xrightarrow{vu}a_2\xrightarrow{1}a_2)$, and $\psi^*h=(h_0,h_2,h_2)$. The first naturality square is the square just computed; the second is $h_2\,1=1\,h_2$. Since $\psi\psi=\psi$, the resulting operation on paths is strictly idempotent. This equality is checked on the three objects of $[2]$, rather than on all categories and all pairs of composable arrows.

Endpoint preservation alone has not made such maps rigid. On $[2]$ the three nondecreasing endpoint-preserving endofunctors have object functions $(0,0,2)$, $(0,1,2)$, and $(0,2,2)$. They respectively place the composite in the second edge, leave the path unchanged, and place the composite in the first edge. This also supplies a concrete reminder of the correction to Proposition~1(b): fixing the endpoints does not force an arbitrary endofunctor of a nonreduced interval to be the identity. The rigidity there belongs to isomorphisms, while the present theory deliberately needs the noninvertible maps.
\end{dcomment}

\begin{dcomment}{exposition}{p25}{Three successive categories of types}
Grothendieck is testing which changes of parameter should count as morphisms of paths. The changes of definition in these pages are purposeful, but a notation retained through all three stages can obscure which fullness assertion remains true. We distinguish the stages here, without altering the source's notation. All three categories have the finite sign words as objects, including the empty word $e$.

Let $T_{\mathrm{all}}$ have every functor $I_\tau\to I_\upsilon$ as a morphism; let $T_{\partial}$ retain just the functors preserving the two chosen endpoints; and let $T_{\uparrow}$ retain those which also preserve the external, linear order of vertices. Thus
\[
 T_{\uparrow}\subseteq T_{\partial}\subseteq T_{\mathrm{all}}.
\]
The category eventually denoted by $\uTy$ in these pages is $T_{\uparrow}$. On page~25 it is initially $T_{\mathrm{all}}$, while pages~29--31 pass through $T_{\partial}$. For $T_{\mathrm{all}}$, realization in $\Cat$ is fully faithful by definition. For $T_{\partial}$, realization is fully faithful into $\varepsilon\backslash\Cat$, where $\varepsilon$ is the discrete category on two objects; forgetting $\varepsilon$ loses fullness. For $T_{\uparrow}$ even the embedding into $\varepsilon\backslash\Cat$ is generally not full. The footnotes retracting fullness on page~26 and the assertion of fullness on page~31 are therefore compatible when read at their respective stages.

Here is a small witness to the last distinction. From the alternating interval $I_{+-+}$ to $I_+$ there is a functor with object function
\[
 (0,1,2,3)\longmapsto(0,1,0,1).
\]
Every forward transition goes to $0\to1$ and the backward transition, whose categorical direction is from its right vertex to its left vertex, also goes to $0\to1$. Both endpoints are preserved, but the vertex function is not nondecreasing. Thus it belongs to $T_{\partial}$ and not to $T_{\uparrow}$. The monotonicity condition excludes this backtracking.

The sign words themselves form a free monoid on $+$ and $-$. Arbitrary adjacent repetitions are allowed; the alternating feature is the alternation of signs between \emph{maximal constant runs}. A nonempty word has a unique expression $\varepsilon_1^{l_1}\cdots\varepsilon_r^{l_r}$ with $l_i>0$ and $\varepsilon_{i+1}=-\varepsilon_i$. Its length is $\sum l_i$, its useful length is $r$, and its critical intervals are these runs with their endpoint vertices restored. The empty word has length and useful length zero and no initial sign. This last exception should remain in force whenever the text speaks of the first generator.

In the undercategory description, $\varepsilon\to I_e$ identifies the two objects of $\varepsilon$. It is a structure map, not an embedding; the hooked arrows drawn on page~29 should not be read as asserting injectivity in this degenerate case. Keeping $e$ is indispensable: it will be the unit for concatenation.
\end{dcomment}

\begin{dcomment}{context}{p27}{The split fibration of all paths}
Grothendieck now uses the machinery of descent to assemble paths of different shapes. The construction is the one in \citet[Expos\'e~VI, \S\S8--9]{GrothendieckSGA1}: a strict contravariant functor to categories supplies a category with a chosen, strictly composable family of Cartesian lifts. The commentary on Chapter~VI, manuscript pages~26 and~53--56, discussed splittings and their universal constructions. Here the splitting needs no choices, because it is given by literal precomposition.

For any of the three shape categories $T$, put $F_X(\tau)=X^{I_\tau}$. An object of its Grothendieck construction is $(\tau,c)$, and a morphism
\[
 (\tau',c')\longrightarrow(\tau,c)
\]
is a pair $(\varphi,v)$ with $\varphi:\tau'\to\tau$ and $v:c'\Rightarrow c\varphi$. If $(\varphi',v'):(\tau'',c'')\to(\tau',c')$, their composite is
\[
 (\varphi,v)(\varphi',v')=
 \bigl(\varphi\varphi',(v\varphi')\circ v'\bigr).
\]
At $i\in I_{\tau''}$ the second component is $v_{\varphi'(i)}v'_i$. Threefold composition is associative because both parenthesizations have component $v_{\varphi'\varphi''(i)}v'_{\varphi''(i)}v''_i$ and the same shape map. The identity is $(1_\tau,1_c)$. These checks explain the word ``\emph{\'evidente}'' in the source without suppressing the variance.

The chosen lift of $\varphi:\tau'\to\tau$ at $(\tau,c)$ is
\[
 (\varphi,1_{c\varphi}):(\tau',c\varphi)\longrightarrow(\tau,c).
\]
To verify its Cartesian property, let a map from $(\omega,d)$ to $(\tau,c)$ have shape $\varphi\theta$ and transformation $w:d\Rightarrow c\varphi\theta$. There is exactly one factor over $\theta$ through the displayed lift, namely $(\theta,w)$. The chosen lifts compose on the nose, since $(c\varphi)\theta=c(\varphi\theta)$. More generally $(\varphi,v)$ is Cartesian exactly when $v$ is invertible: it factors as the vertical arrow $(1,v)$ followed by this chosen lift, and a vertical Cartesian arrow is an isomorphism.

For $T=T_{\partial}$ or $T_{\uparrow}$ the endpoint functors send $(\varphi,v)$ to $v_{0_{\tau'}}$ and $v_{1_{\tau'}}$. Their targets are $c(0_\tau)$ and $c(1_\tau)$ because $\varphi$ preserves endpoints. The formula for composition immediately gives functoriality. Likewise a functor $H:X\to Y$ acts by $(\tau,c)\mapsto(\tau,Hc)$ and $(\varphi,v)\mapsto(\varphi,Hv)$, preserving the chosen Cartesian lifts. Restricting the permitted shape maps is literally base change of this split fibration, which proves the fiber-product assertion on page~30.

Size causes no mathematical obstruction here. Fix universes $\mathcal U\in\mathcal V$. The set of finite words and their finite interval functors is $\mathcal U$-small. If $X$ is a $\mathcal V$-small category, each $X^{I_\tau}$ and the total construction are $\mathcal V$-small. One may therefore take $X=\Cat_{\mathcal U}$ after working in a sufficiently larger universe. This supplies the convention which the source deliberately leaves implicit.
\end{dcomment}

\begin{dcomment}{examples}{p28}{Changing shape need not have a direct image}
The projection to the category of types is a fibration because paths can always be restricted. Extending a path to a new shape asks a different question, and generally requires colimits in its coefficient category. This is the usual distinction between a fibration and a bifibration; compare \citet[Expos\'e~VI, \S10 and \S11(f)]{GrothendieckSGA1}.

Consider the interval $J=I_{-+}=(0\leftarrow1\to2)$ and its unique collapse $q:J\to e$. This collapse belongs even to $T_{\uparrow}$. Its inverse-image functor is the diagonal
\[
 q^*:X\longrightarrow X^J,
\]
sending an object to its constant span. If an opCartesian lift over $q$ existed at every span $c$, its target $q_!c$ would represent the functor
\[
 x\longmapsto\operatorname{Nat}(c,\mathrm{const}_x).
\]
Indeed every map from $c$ over $q$ to a constant path would factor uniquely through the opCartesian lift. This is precisely the universal property of the pushout of the span $c(0)\leftarrow c(1)\to c(2)$.

Take $X=J$ itself and $c=1_J$. There is no cocone on this diagram: neither $0$ nor $2$ maps to the other, and neither maps to $1$. Thus $q_!c$ does not exist, and the path projection is not an opfibration in this example. This does not reveal a gap in Grothendieck's assertion, which is only that the construction is split fibrational. It identifies precisely what would be added by a stronger assertion.

For $X=\mathrm{Set}$ the same calculation has a complete answer. Let $A\leftarrow C\to B$ be a span. Then $q_!c$ is the quotient of $A\amalg B$ by the equivalence relation generated by identifying the two images of each $z\in C$. A map from this quotient to $S$ is exactly a pair $A\to S$, $B\to S$ which agrees on $C$, so it is exactly a natural transformation $c\Rightarrow\mathrm{const}_S$. For a general change of shape $\varphi$ and a cocomplete $X$, left Kan extension supplies the corresponding left adjoint to $\varphi^*$. Such extra adjoints come from $X$; the split fibration itself is available without them.
\end{dcomment}

\begin{dcomment}{development}{p29}{How much endpoint preservation is forced}
The correction on page~29 gives exactly what is needed for endpoint maps obtained by taking components of a natural transformation. There is a slightly larger possibility if one allows those components to be followed by canonical arrows inside the target interval. Describing it clarifies why endpoint preservation is a structural choice, and why concatenation will justify the stricter choice.

\textbf{Endpoint evaluation criterion.}\ours{}
Let $T$ be a subcategory of $T_{\mathrm{all}}$ with all sign words as objects. Its path construction admits functors $\sigma_X,\beta_X:\int F_X\to X$, natural in $X$, restricting in each fiber to evaluation at the specified source and target, if and only if every $\varphi:\tau'\to\tau$ in $T$ admits arrows
\[
 a_\varphi:\varphi(0_{\tau'})\longrightarrow0_\tau,
 \qquad
 b_\varphi:\varphi(1_{\tau'})\longrightarrow1_\tau
 \quad\text{in }I_\tau.
\]
When they exist, these endpoint functors are unique. The permitted functors form a category $T_{\mathrm{ev}}$ containing $T_{\partial}$, and it is the largest such subcategory of $T_{\mathrm{all}}$.

\emph{Proof.} For necessity take $X=I_\tau$ and the universal path $c=1_{I_\tau}$. The Cartesian arrow
\[
 (\varphi,1_\varphi):(\tau',\varphi)\longrightarrow(\tau,1_{I_\tau})
\]
must be sent by $\sigma_X$ to an arrow $\varphi(0_{\tau'})\to0_\tau$, and by $\beta_X$ to an arrow $\varphi(1_{\tau'})\to1_\tau$. Each is unique when it exists, because an interval is thin. For composable maps $\tau''\xrightarrow{\psi}\tau'\xrightarrow{\varphi}\tau$, the arrows
\[
 a_\varphi\circ\varphi(a_\psi),\qquad
 b_\varphi\circ\varphi(b_\psi)
\]
have precisely the endpoints required for $\varphi\psi$. Identities satisfy the requirement as well. Hence the permitted functors form a category.

To construct the endpoint maps, define on $(\varphi,v):c'\to c$
\[
 \sigma_X(\varphi,v)=c(a_\varphi)\circ v_{0_{\tau'}},
 \qquad
 \beta_X(\varphi,v)=c(b_\varphi)\circ v_{1_{\tau'}}.
\]
For a second morphism $(\psi,w):c''\to c'$, naturality of $v$ at $a_\psi$ gives
\[
 v_{0_{\tau'}}c'(a_\psi)
   =c\varphi(a_\psi)v_{\psi(0_{\tau''})}.
\]
After composing with $c(a_\varphi)$ and $w_{0_{\tau''}}$, this says that the product of the two source arrows is
\[
 c(a_{\varphi\psi})
 v_{\psi(0_{\tau''})}w_{0_{\tau''}},
\]
which is the source arrow assigned to the composite. The same argument with $b_\psi$ proves the target identity. Identity morphisms and naturality in $X$ follow directly from the formulas.

Finally, every $(\varphi,v)$ factors as a vertical morphism $(1,v)$ followed by a chosen Cartesian arrow. The endpoint functors are prescribed on vertical arrows. Naturality in $X$ applies to $c:I_\tau\to X$ and forces their values on the Cartesian arrow to be $c(a_\varphi)$ and $c(b_\varphi)$. Thus the formulas are forced, proving uniqueness and maximality.\hfill$\square$

For example, the functor $e\to I_+$ selecting $0$ belongs to $T_{\mathrm{ev}}$: its source comparison is an identity and its target comparison is $0\to1$. It does not preserve both endpoints. The functor selecting $1$ does not belong to $T_{\mathrm{ev}}$, since there is no arrow $1\to0$. This already prevents endpoint evaluation on $T_{\mathrm{all}}$.

The enlargement does not solve the concatenation problem. Concatenate two copies of the admissible map $e\to I_+$ selecting $0$. The common vertex of $e*e=e$ would have to map both to vertex $0$ of the first target interval and to vertex $1$, the starting vertex of the second interval in $I_{++}$. The two prescribed maps therefore do not glue. Strict endpoint preservation makes precisely these images agree. The theorem is an elementary specialization of the general description of functors out of a split fibration in \citet[Expos\'e~VI, Proposition~12.1]{GrothendieckSGA1}. We checked that description and the Yoneda treatment cited above; we did not find this particular maximal class of interval maps stated there. The marker records this explicit deduction and proof, without a claim of priority.
\end{dcomment}

\begin{dcomment}{status}{p30}{The fraction test and hammock localization}
Grothendieck's two tests concern different levels of the construction. Concatenation should be an operation on actual paths and their morphisms; localization should arise after paths are subjected to specified identifications or homotopical comparisons. The category $\uCh(X)$ constructed here has not yet been asserted to satisfy the universal property of a localization.

For a category $\mathcal M$ and a specified class $\Sigma$ of arrows, ordinary localization is already constructed by \citet[Chapter~I, \S1.1 and Lemma~1.2]{GabrielZisman1967}. Its arrows are composable words of forward arrows and formal inverses of arrows in $\Sigma$, modulo the category relations and inverse relations. A functor out of the quotient is consequently equivalent to a functor from $\mathcal M$ sending $\Sigma$ to isomorphisms. This existence statement does not require the Ore hypotheses needed for a calculus with uniformly short fractions.

The homotopical comparison already available before these notes is quite specific. In \citet[Definition~2.1 and Propositions~2.2 and~3.1]{DwyerKan1980Calculating}, a hammock is a finite grid of zigzags, with backward horizontal arrows and vertical arrows required to be weak equivalences. Identity columns are removed and adjacent columns pointing in the same direction are composed. Their mapping simplicial sets are equivalent to the standard simplicial localization, and their sets of components recover ordinary localization:
\[
 \pi_0L^H(\mathcal M,\Sigma)(x,y)
   \cong\Hom_{\mathcal M[\Sigma^{-1}]}(x,y).
\]
Thus there is a proven construction in which changing the length and subdivision of zigzags is accompanied by the higher comparison data needed for mapping spaces.

This is a comparison of mathematical mechanisms, not an identification of the constructions. The present $\uCh(X)$ allows arbitrary vertical natural transformations, has not specified which backward arrows may be inverted, and uses a single ordinary category of types. The hammock conditions, the fixed endpoints, and the higher simplices must all be supplied before a localization theorem can be invoked. The source's demand that the fraction theory ``work'' is a design requirement at this point, not a research conjecture left open in these pages. Its exact implementation in Grothendieck's subsequent text should be checked on its own definitions.
\end{dcomment}

\begin{dcomment}{exposition}{p32}{A functor cannot cross a turn in one run}
Grothendieck begins the structure theorem by cutting an interval where its direction changes. Each piece is a chain, and the image of a chain remains comparable even when its length changes. This elementary observation is what allows a map of long zigzags to be described by a finite amount of ordered data.

For a nonempty word $\tau$ let
\[
 0=s_0<s_1<\cdots<s_r=n
\]
be its critical vertices, and let $K_i=[s_{i-1},s_i]$ be its maximal constant-sign runs, regarded as full subcategories. The category $I_\tau$ is obtained by gluing the $K_i$ at their common endpoint vertices. To see the universal property directly, choose functors $f_i:K_i\to Y$ agreeing at shared endpoints. They determine the image of every generating transition of $I_\tau$. Each nonidentity arrow of $I_\tau$ is the unique composite of the transitions along a single run, so its image is forced to be the corresponding composite in $Y$. The only composition equations are associativity along those runs, and these hold in $Y$. Thus the $f_i$ extend uniquely to a functor. For $Y=I_{\upsilon}$, imposing the values of $f_1(s_0)$ and $f_r(s_r)$ additionally imposes endpoint preservation, exactly as on page~32.

Now let $K$ be a critical run of the source and suppose $f|K$ is nonconstant. The image vertices are pairwise comparable in $I_{\upsilon}$: any two source vertices are joined by an arrow in one direction, whose image joins their images. Let $a<b$ be the smallest and largest distinct image vertices in the external order. Their comparability means that all target transitions between $a$ and $b$ have one sign. They therefore lie in a unique maximal run $J'$, and so do all the image vertices between them. This proves both the existence and uniqueness in Proposition~7.

If the source run has sign $\varepsilon$ and $J'$ has sign $\eta$, the object function of $f|K$ is nondecreasing when $\varepsilon=\eta$ and nonincreasing when $\varepsilon=-\eta$. A nonconstant function which is nondecreasing in the external order cannot be of both sorts. Hence an increasing functor sends every nonconstant source run into a target run of the same sign. This proves the sign assertion as well; no endpoint hypothesis has been used.

The preceding paragraph of the source overstates the constant case. A vertex strictly inside a long target run belongs to one critical interval; an interior \emph{critical} vertex belongs to two; an endpoint of a positive-length interval belongs to one. For instance, the constant map with value $1$ in $I_{++}=[2]$ factors through its sole critical interval, although $1$ is not an endpoint. This is a printed slip, retained in the text, and it does not affect Proposition~7 because that proposition assumes nonconstancy. There is also a harmless zero-length exception to the footnote's comparable-subset lemma: $I_e$ has one vertex and no critical runs. A nonconstant map cannot have this target, so the proposition itself remains valid without an added exception.
\end{dcomment}

\begin{dcomment}{exposition}{p34:sous-intervalle réduit qu'il définit}{Critical vertices are crossed, not necessarily preserved}
The next step records which turns in the target must be encountered by a map. This is weaker than preservation of all critical vertices, and the distinction is needed when the target is not reduced. It also prevents an ordinary interval of integers from being confused with a connected full subcategory of an interval category.

Let $s_0,\ldots,s_r$ be the critical vertices of the source and put $t_i=f(s_i)$. Consecutive $t_i$ are comparable, since they are images of the endpoints of a source run. Thus the full subcategory on the \emph{set} $\{t_0,\ldots,t_r\}$ is connected. Suppose $t_i$ and $t_j$ lie on opposite sides of an interior critical vertex $z$ of the target. If no $t_k$, $i\leq k\leq j$, were equal to $z$, some consecutive pair in this finite sequence would lie strictly on opposite sides of $z$. Such a pair cannot be comparable: the two transitions adjoining an interior critical vertex have opposite signs. This contradiction proves that $z=t_k$ for some $k$ in the indicated range. The argument permits backtracking and repetitions; it does not require monotonicity.

If $f$ preserves the two endpoints, take $i=0$ and $j=r$. Every target critical vertex is consequently the image of a source critical vertex. This is the content of the corollary, whose first occurrence of $I_{\tau'}$ should be read as $I_\tau$ when it describes the $s_i$. The stronger inclusion in the other direction is false. In the displayed example $I_{+-+}\to I_{++}$, the object function is
\[
 (0,1,2,3)\longmapsto(0,1,1,2).
\]
All four source vertices are critical. Their image contains the two target endpoints and also its noncritical vertex $1$; the middle source transition collapses. The example therefore gives exactly the one-sided containment asserted in the text.

There is a second printed overstatement just before the corollary. Even for increasing $f$, the full subcategory on the $t_i$ need not have \emph{every} integer vertex between $t_0$ and $t_r$. Take the identity of $I_{++}$. Its source critical vertices are $0,2$, so their image is $\{0,2\}$, a full subcategory isomorphic to $[1]$, whereas the ambient interval $[0,2]$ has three vertices. The correct statement is that the image subcategory has least and greatest vertices $t_0,t_r$, its order-convex hull is $[t_0,t_r]$, and it contains every ambient critical vertex in that hull. This is all that the argument uses.

One can now see why the image of $\Cr(\tau,f)\to\Cr(\upsilon)$ is an interval in the ordered set of target runs. Join successive numbers $t_i$ by real line segments, obtaining a continuous walk in $[0,m]$. Each nonconstant segment is contained in the unique target run attached to $K_i$. If the walk reaches two distinct target runs, it must traverse every intervening run, and each such traversal occurs on at least one nonconstant segment. Conversely, a target run not encountered by the walk cannot be attached to any $K_i$. The same argument shows that every target run meeting the segment between $f(0)$ and $f(n)$ in a positive-length segment is encountered. In particular, endpoint preservation makes the run map surjective. This completes the proofs of Corollary~2(a)--(c), including the parts for maps which backtrack.
\end{dcomment}

\begin{dcomment}{exposition}{p36}{The length budget over a target run}
Monotonicity turns the preceding coverage statement into an additive formula. A map may pause, but it cannot spend the same portion of the target twice. This is the numerical mechanism behind Corollaries~3--5.

Let $f:I_\tau\to I_\upsilon$ be increasing and preserve endpoints. For a source critical run $K_i=[s_{i-1},s_i]$, write
\[
 d_i=f(s_i)-f(s_{i-1})\geq0.
\]
When $d_i>0$, Proposition~7 associates to $K_i$ the unique target run containing its image; this is $\Cr(f)(K_i)$. Fix a target run $J'=[a,b]$. The nonconstant segments $[f(s_{i-1}),f(s_i)]$ attached to $J'$ cover $[a,b]$, meet only at consecutive endpoints, and are encountered in their natural order. Coverage follows because $f(0)=0$, $f(n)=m$ and every intervening critical vertex is attained; disjointness of interiors follows from the monotonicity of the sequence $f(s_i)$. Their lengths therefore add to $b-a$:
\[
 b-a=\sum_{\substack{K_i\in\Cr(\tau,f)\\\Cr(f)(K_i)=J'}}d_i.
\]
This is the formula for $\pi_{\mathrm n}(K_i,f)$ in the source. It is valid even if $f$ skips noncritical target vertices, because length measures the whole segment spanned by the image, not the number of image vertices.

Every summand displayed is a positive integer. Consequently, if $k_{J'}=|\Cr(f)^{-1}(J')|$, then
\[
 1\leq k_{J'}\leq b-a,
\]
and equality on the right holds precisely when every positive $d_i$ in the sum is $1$. The printed arrow in Corollary~3 is a typesetting slip for $\leq$, and its unprimed $J$ in the singleton should be $J'$. The subsequent discussion of equality gives the intended reading independently of our proof.

For a first calculation take $I_{+-+}\to I_{++}$ with vertex function $(0,1,1,2)$. Its three source runs have lengths of image $1,0,1$; the single target run has length $2$. Thus $k_{J'}=2$ and the upper bound is attained. For the identity $I_{++}\to I_{++}$, the sole source run has image length $2$, so $k_{J'}=1<2$. Both maps are increasing and preserve endpoints. What is counted is the number of nonconstant \emph{runs}, not the number of nonidentity generating arrows.

The alternative summation on page~36 may harmlessly include every global source run contained in $f^{-1}(J')$: the added terms have zero length. Its parenthetical identification of these runs with all critical runs of the preimage is, however, too literal. A preimage can contain a truncated piece of a neighboring global run on which $f$ is constant. For example, take
\[
 f:I_{++--}\longrightarrow I_{+-},\qquad
 (0,1,2,3,4)\longmapsto(0,1,1,1,2).
\]
For $J'=[0,1]$, its preimage has vertices $0,1,2,3$ and critical runs $[0,2]$ and $[2,3]$. Only $[0,2]$ is a full global source run contained in the preimage; $[2,3]$ is a proper part of the global run $[2,4]$. The extra piece is constant with value $1$, hence contributes zero. The equality of sums survives, but equality of the indexing sets does not. One should either keep the global-run indexing or explicitly allow such zero-contribution boundary pieces.
\end{dcomment}

\begin{dcomment}{examples}{p38}{A reduced target remembers the surviving transitions}
Grothendieck next specializes the length formula to a target with no redundant vertices. Every target run then has length one. There is exactly one nonconstant source critical run over it, and that source run can cross its two target vertices only once. This converts a numerical constraint into the combinatorial invariant used in Proposition~8.

More explicitly, on a source run $K=[s,t]$ mapping onto the two vertices $a,a+1$, the nondecreasing vertex function has the form
\[
 f(s)=\cdots=f(i-1)=a,
 \qquad
 f(i)=\cdots=f(t)=a+1
\]
for exactly one $i$ with $s<i\leq t$. Hence exactly one transition in $K$ has nonidentity image. The run and its target have the same sign, so this transition maps to the target's unique transition in its categorical direction, including when that direction is backward. Performing this construction for all target runs gives the bijection of Corollary~5. It is increasing because the source runs and their crossings occur in the same external order as the target runs. The triangle with $\alpha$ commutes because $\alpha$ forgets which transition was selected and remembers its containing source run.

Take $\tau=++--++$ and $\upsilon=+-+$. Choose transitions $2,4,6$ in the source. The associated vertex map is
\[
 (0,1,2,3,4,5,6)\longmapsto(0,0,1,1,2,2,3).
\]
The first source run maps across the first forward target edge at transition~2; the second maps across the backward target edge at transition~4; the third crosses the final forward edge at transition~6. The other three transitions become identities. If a path of type $\upsilon$ in $X$ is
\[
 x_0\xrightarrow{a}x_1\xleftarrow{b}x_2\xrightarrow{c}x_3,
\]
restriction along this map is the path
\[
 x_0\xrightarrow{1}x_0\xrightarrow{a}x_1
 \xleftarrow{1}x_1\xleftarrow{b}x_2
 \xrightarrow{1}x_2\xrightarrow{c}x_3.
\]
In each of the three source runs either of its two transitions could instead be the surviving one. These choices give $2^3=8$ different maps to this reduced target, all explicitly determined by their three crossing positions. The general theorem on the next pages will allow arbitrary sign words and express this determination as a signed-subsequence bijection.

When the target is $e$, there are no target runs and no nondegenerate transitions. The unique endpoint-preserving map $I_\tau\to I_e$ collapses every transition. Both sides of the bijection in Corollary~5 are empty, so the assertion includes this case. When the source is $e$ and the target has positive length, no endpoint-preserving map exists. These two degenerate cases explain why the later subset description needs the empty word as well as the alternating nonempty words.
\end{dcomment}

\begin{dcomment}{exposition}{p40:se donner une suite strictement croissante}{Choosing the transitions that survive}
The reduced target allows Grothendieck to replace a functor by a selection of edges. This is the first form of the normal form that will make concatenation manageable: each edge of the target has one identifiable ancestor in the source. The hypothesis that the target is reduced is doing exactly the work needed here.

Write the source signature as $\tau=(\varepsilon_1,\ldots,\varepsilon_n)$ and the reduced target as $\tau'=(\varepsilon'_1,\ldots,\varepsilon'_m)$. An increasing endpoint-preserving functor has a nondecreasing object map
\[
 f:[0,n]\longrightarrow[0,m],\qquad f(0)=0,\quad f(n)=m.
\]
The image of the $i$th transition is either an identity or a nonidentity arrow joining $f(i-1)$ to $f(i)$. In a reduced interval every nonidentity arrow is a transition: two successive transitions have opposite signs and cannot compose. Consequently
\[
 f(i)-f(i-1)\in\{0,1\}.
\]
Exactly $m$ differences are $1$. If their indices are $i_1<\cdots<i_m$, preservation of arrow direction gives $\varepsilon_{i_j}=\varepsilon'_j$. Conversely such a signed subsequence defines
\[
 f(k)=\#\{j\mid i_j\leq k\}.
\]
On a selected transition this map uses the corresponding target transition; on every other transition it uses an identity. These choices extend uniquely to a functor because the source is the free category on its oriented linear graph. They preserve the two endpoints and recover the selected subsequence. This proves Proposition~8, including its empty-target case: the empty subsequence gives the unique functor to the point.

For example, from the word $++-+$ to the reduced word $+-+$ there are exactly two morphisms. The negative transition must be the third source edge, the last positive transition must be the fourth, and the first positive transition may be either the first or the second. Their object maps are respectively
\[
 (0,1,1,2,3),\qquad (0,0,1,2,3).
\]
The two paths retain different occurrences of the first positive step. Thus the type records positions, not merely the numbers of positive and negative transitions.

For a single-edge target the same argument gives the stated representability of positive and negative transitions. A positive edge at position $i$ determines the map to $\Delta_+$ which is $0$ before that edge and $1$ after it; a negative edge determines the analogous map to $\Delta_-$. Naturality remains to be explained because an edge can be stretched over several target edges. The next construction supplies exactly that explanation. When the target is not reduced, the selection alone is insufficient: the unique morphism $\Delta_+\to\Delta_+^2$ stretches one edge across two, so the size of each surviving image must also be retained.
\end{dcomment}

\begin{dcomment}{examples}{p42}{A stretched edge and the containment convention}
Grothendieck is assigning every target transition to the unique source transition whose image crosses it. One small example fixes both the meaning of this assignment and two printed slips immediately before Corollary~2.

Let $f:\Delta_+\to\Delta_+^2$ have object map $f(0)=0$, $f(1)=2$. There is one surviving source transition, so $N=1$, $t_0=0$ and $t_1=2$. The two target transitions have underlying vertex segments $[0,1]$ and $[1,2]$. Both are contained in $[t_0,t_1]=[0,2]$, and neither equals it. Thus the printed relation
\[
 [i-1,i]=[t_{j-1},t_j]
\]
must be read mathematically as $[i-1,i]\subseteq[t_{j-1},t_j]$. The parenthesis identifying the part of $f$ on $v_j$ as $[j-1,j]$ has the same problem: that part is $[t_{j-1},t_j]$. Both readings occur in the edition; they are retained in the source and repaired here. Corollary~2 itself uses containment and is correct.

The intervals $[t_{j-1},t_j]$ may meet at a vertex, but their sets of transitions are disjoint. Since $t_0=0$, $t_N=m$ and the $t_j$ are strictly increasing, these transition sets partition $\{1,\ldots,m\}$. A transition cannot lie in two of them: its two endpoints would then have to belong to their intersection, which contains at most one vertex. This proves uniqueness without assuming that the target is reduced. Functoriality forces every edge in the block assigned to $v_j$ to have the sign of $v_j$.

The words ``source'' and ``but'' in the preceding display concerning the $v_j$ must here be read as the first and last vertices in the traversal order. On a retrograde transition they are the reverse of its categorical domain and codomain. Writing $f(i_j-1)$ and $f(i_j)$ avoids this ambiguity. If the target is the point there are no target transitions and the pullback map of transition sets is the unique map from the empty set; no first or last surviving transition should be invoked.
\end{dcomment}

\begin{dcomment}{exposition}{p43:Ainsi, le corollaire 1 plus haut}{Edges reverse the direction of a map}
The contravariant transition map contains more information than its notation initially suggests. It determines the whole functor, and gives a particularly economical modern description of the final category of types. This observation explains why the later calculations require only ordered lists and nonnegative integers.

Let $\mathcal O_{\pm}$ have as objects finite linearly ordered sets equipped with a sign at each element, including the empty set, and as arrows the nondecreasing maps which preserve signs. We use the skeletal version with underlying sets $[1,n]$. For $f:\tau\to\tau'$ in $\uTy$, where $\tau$ and $\tau'$ have lengths $n$ and $m$, define
\[
 q_f:[1,m]\longrightarrow[1,n],\qquad
 q_f(j)=\min\{i\mid j\leq f(i)\}.
\]
The endpoint conditions ensure that this minimum exists when $m>0$; its defining inequalities are
\[
 f(q_f(j)-1)<j\leq f(q_f(j)).
\]
This is Grothendieck's $\Tr(f)^*$. It is nondecreasing, and it preserves signs because the image of a source transition is a monochromatic string of target transitions of that same sign.

Conversely, suppose $q:[1,m]\to[1,n]$ is nondecreasing and sign-preserving. Put
\[
 F_q(i)=\#q^{-1}([1,i]),\qquad 0\leq i\leq n.
\]
The fibre $q^{-1}(i)$ is a consecutive block, possibly empty, whose target signs are all $\varepsilon_i$. Hence the consecutive vertices $F_q(i-1)$ and $F_q(i)$ are joined in the target by an arrow of the required direction, or coincide when the fibre is empty. Sending the $i$th source transition to this arrow defines a unique functor. Its object map is nondecreasing, begins at $0$, and ends at $m$. The inequalities above show $q_{F_q}=q$, while counting the target edges crossed by the first $i$ source edges gives $F_{q_f}(i)=f(i)$. These constructions are inverse, including the empty cases. There is no map from the point to a nonpoint interval in $\uTy$, just as there is no function from a nonempty edge set to the empty one.

If $g:\tau'\to\tau''$, a target edge of $\tau''$ has a unique ancestor in $\tau'$ and then a unique ancestor in $\tau$. Equivalently, applying the displayed inequalities twice gives
\[
 q_{gf}=q_fq_g.
\]
Identities induce identities. We therefore have an isomorphism of skeletal categories
\[
 \uTy\ \cong\ \mathcal O_{\pm}^{\mathrm{op}}.
\]
The reversal of variance is essential: a functor stretches a source edge into a target block, whereas $q_f$ sends each member of that block back to its source edge. The signs can be forgotten for the construction on ordered sets, as Grothendieck observes, but must be restored to say when the resulting vertex map is a functor between the oriented interval categories.

This identification also establishes the naturality promised in Corollary~1. If an edge of sign $\varepsilon$ represents $h:\tau'\to\Delta_\varepsilon$, the edge representing $hf$ is $q_f$ of the edge representing $h$. It is precisely ordinary precomposition in the representable presheaf $\Hom_{\uTy}(-,\Delta_\varepsilon)$. The underlying single-sign calculation is the familiar ordinal description of the simplicial category; see \citet[Appendix~D.1, pp.~713--715]{AguiarMahajan2010}. The signed formulation follows by retaining the labels throughout the inverse constructions above.
\end{dcomment}

\begin{dcomment}{examples}{p44}{Weight vectors and composition by block sums}
The outgoing classification now has coordinates which are useful for actual calculations. For a fixed source $\tau=(\varepsilon_1,\ldots,\varepsilon_n)$, put
\[
 a_i=f(i)-f(i-1)=\#q_f^{-1}(i),\qquad
 A_i=\sum_{h=1}^{i}a_h,\quad A_0=0.
\]
Every $a_i$ is a nonnegative integer. The target is the word
\[
 \tau'=\varepsilon_1^{a_1}\varepsilon_2^{a_2}\cdots\varepsilon_n^{a_n},
 \qquad f(i)=A_i,
\]
where exponent zero means omission. These equations prove the claimed unrestricted parametrization of $\Ob(\tau\backslash\uTy)$: they construct an outgoing morphism from every weight vector and recover the vector from every outgoing morphism. When adjacent source signs agree, their target blocks may merge into one critical interval, but their individual lengths remain part of the morphism.

Take the source $+-++-$ and the weights $(2,0,1,3,0)$. Its target is $++++++$, and its vertex map is
\[
 (0,2,2,3,6,6).
\]
The pullback on target transitions is $(1,1,3,4,4,4)$. Both negative source edges have been collapsed. The two consecutive positive source edges at positions $3$ and $4$ survive separately, even though their target images lie in the same positive run.

To see composition, give the six target edges the further weights
\[
 (b_1,\ldots,b_6)=(1,0,2,1,0,3).
\]
The second map has vertex values $(0,1,1,3,4,4,7)$ and target $+^7$. The composite has weights $(1,0,2,4,0)$, since its $i$th weight is
\[
 c_i=\sum_{j=A_{i-1}+1}^{A_i}b_j.
\]
Thus its vertex map is $(0,1,1,3,7,7)$, which is also obtained by direct substitution into the two displayed object maps. An empty block contributes zero; a block of length greater than one contributes a sum. This is the numerical content of composition in Section~6.

The formula gives a useful direct verification of associativity. Given a third weight vector, first replacing every $b_j$ by the sum over its descendants and then summing over the $a_i$-block is the same finite sum as first merging the $b_j$-blocks and then summing over the merged block. The order and boundaries of the summands are unchanged. No cancellation, choice of representative, or equivalence relation enters this calculation.
\end{dcomment}

\begin{dcomment}{exposition}{p45}{Fibres are components of the collapsed-edge graph}
Grothendieck wants a geometric reconstruction of the same weight vector. The intended reconstruction is correct once the graph of collapsed transitions is kept, but the printed description in terms of a set of vertices loses information that is needed to find its components.

Consider $\tau=+++$ with weights $(0,1,0)$. The corresponding map to $+$ is
\[
 f(0)=f(1)=0,\qquad f(2)=f(3)=1.
\]
The vertices incident to a collapsed transition form the whole set $\{0,1,2,3\}$. As a subset of the integers this set is connected. Consequently the assertion that $f$ is constant on its integer-connected components would make $f$ constant everywhere, contrary to the displayed functor. This is a printed mathematical slip, not a defective proof of the weight classification.

Instead, take the graph with all source vertices and only those edges $[i-1,i]$ for which $a_i=0$. Its components are consecutive strings of zero-weight edges, together with any isolated vertices. In the example they are $\{0,1\}$ and $\{2,3\}$. Since
\[
 f(j)-f(i)=\sum_{h=i+1}^{j}a_h\qquad(i\leq j),
\]
nonnegativity shows that two vertices have the same image if and only if every intervening edge has weight zero. The graph components are therefore exactly the nonempty vertex fibres of $f$. On each surviving transition the endpoint values are determined by the cumulative sum of the earlier weights. This proves both existence and uniqueness of the proposed geometric reconstruction after the correction.

The same observation produces a canonical factorization. First collapse the zero-weight edges, leaving the signed subword consisting of the active source edges; then stretch its successive edges by their positive weights. The first vertex map is surjective, the second strictly increasing, and both preserve the endpoints. Uniqueness follows either from the fibre components or from the unique image factorization of $q_f$ in the category of signed ordered sets. This factorization will let us distinguish a combinatorial contraction from a contraction which has an inverse on one side.
\end{dcomment}

\begin{dcomment}{exposition}{p47}{Reading the morphism from its target}
The incoming classification organizes a morphism by the places where its nonconstant pieces meet. It is the same normal form viewed from the target: choose the visited vertices, and at each visited vertex record the entire interval that has been collapsed there. This description is useful because the intervening noncollapsed transitions are then forced.

Let $f:I_\tau\to I_{\tau'}$ be increasing and endpoint-preserving, with target length $m$. The uniform definition is
\[
 E'=\operatorname{im}(f)\subseteq[0,m].
\]
When $m>0$, it agrees with the printed definition by endpoints of nondegenerate images. Indeed, take a vertex value $t$ assumed by $f$. If not every source vertex has value $t$, connectedness of the underlying source path supplies an edge from the nonempty set $f^{-1}(t)$ to its complement. That edge is nondegenerate and one of its image endpoints is $t$. Conversely every such endpoint is a value of $f$. This argument does not require monotonicity, and supplies the elementary lemma requested in footnote~13, provided the functor is nonconstant.

For $m=0$ all transition images are identities, so the printed endpoint definition would instead give $E'=\varnothing$. The later claim that $E'$ contains the endpoint, and the definition $N=\#E'-1$, then fail. Taking $E'=\operatorname{im}(f)=\{0\}$ repairs the nullary case and agrees with the explicit description of $N=0$ on manuscript page~47. This is an omitted boundary convention; it does not change the nonconstant construction.

Write the distinct elements of $E'$ as
\[
 0=t_0<t_1<\cdots<t_N=m.
\]
Every two consecutive values are joined by a target arrow: the source transition between their two adjacent nonempty fibres has exactly these image endpoints. Hence the segment from $t_{j-1}$ to $t_j$ is monochromatic, with sign $\varepsilon_j$. Equivalently, $E'$ contains the intrinsic critical vertices of the target, namely its endpoints and every place where the signs change. The notation $\mathrm{Cr}(\tau')$ has changed meaning at this point in the manuscript: here it denotes these vertices, whereas earlier it denoted critical intervals. Keeping these two sets distinct removes a potential confusion in the formula.

The fibre over $t_j$ is a nonempty consecutive interval of vertices; let its signature be $\tau_j$ and its length $\lambda_j$. Reading the source from left to right gives exactly
\[
 \tau=\tau_0\Delta_{\varepsilon_1}\tau_1
        \Delta_{\varepsilon_2}\cdots
        \Delta_{\varepsilon_N}\tau_N.
\]
Conversely, choose any subset $E'$ containing the critical vertices and any words $\tau_0,\ldots,\tau_N$. Form this source word. Collapse the $\tau_j$-block to $t_j$, and map the transition between consecutive blocks to the target arrow joining $t_{j-1}$ and $t_j$, in the direction prescribed by $\varepsilon_j$. This defines an increasing endpoint-preserving functor; its image is exactly $E'$, and its fibres have the prescribed words. Each arrow is uniquely determined by its endpoints in an interval category, so there is no additional gluing datum. The two constructions are inverse. For a point target they reduce to the arbitrary word $\tau_0$ and its unique map to the point.

For use in computations, set
\[
 l_j=t_j-t_{j-1},\qquad
 B_j=j+\sum_{h<j}\lambda_h.
\]
The correct full formula is
\[
 f(k)=t_j=\sum_{h=1}^{j}l_h
 \quad\text{for } B_j\leq k\leq B_j+\lambda_j,
 \qquad 0\leq j\leq N.
\]
The source's explicit table gives $l_2$ on the third fibre, where the value must be $l_1+l_2$; its final row similarly needs the corresponding cumulative sum. The formula above repairs all rows at once. For example, take target $++-$, visited vertices $E'=\{0,2,3\}$, and fibre words $\tau_0=-$, $\tau_1=+$, $\tau_2$ empty. Then the source is $-++-$ and the object map is $(0,0,2,2,3)$. The first positive transition crosses the two target positive edges, the next is collapsed, and the final negative transition crosses the final target edge. The two jumps have sizes $2$ and $1$, while the final vertex value is their sum $3$.
\end{dcomment}

\begin{dcomment}{development}{p47}{Which contractions admit a section}
The normal form makes it possible to ask precisely when a deleted collection of transitions can be restored functorially. This is useful for the later path formalism: a surjection of underlying vertex sets need not be a split epimorphism in the category of signed types. The signs can obstruct every possible section.

\textbf{Section criterion for a contraction.}\ours{}
Let $\tau=(\varepsilon_1,\ldots,\varepsilon_n)$ and let $E=\{e_1<\cdots<e_k\}$ be a subset of its transition positions. Let
\[
 p_E:\tau\longrightarrow\tau|_E
\]
be the contraction with weights $1$ at the positions in $E$ and $0$ elsewhere. If $k=0$, it has a section exactly when $n=0$. If $k>0$, it has a section exactly under the following conditions. Every sign before $e_1$ equals $\varepsilon_{e_1}$, and every sign after $e_k$ equals $\varepsilon_{e_k}$. In each open gap $(e_j,e_{j+1})$, the sequence of signs is a possibly empty string of $\varepsilon_{e_j}$ followed by a possibly empty string of $\varepsilon_{e_{j+1}}$. When these conditions hold, the number of sections is
\[
 \prod_{\substack{1\leq j<k\\
                  \varepsilon_{e_j}=\varepsilon_{e_{j+1}}}}
       (e_{j+1}-e_j).
\]
An empty product is $1$.

\emph{Proof.} Under the signed-edge anti-isomorphism, $p_E$ corresponds to the inclusion $E\hookrightarrow[1,n]$. A section of $p_E$ therefore corresponds exactly to a nondecreasing sign-preserving retraction
\[
 r:[1,n]\longrightarrow E,\qquad r(e_j)=e_j.
\]
If $E$ is empty such a map exists exactly when its domain is empty. Assume henceforth $k>0$. Monotonicity forces every index before $e_1$ to map to $e_1$, and every index after $e_k$ to map to $e_k$. This proves the two exterior conditions, and these parts of $r$ are unique.

For $e_j<i<e_{j+1}$, monotonicity forces
$e_j\leq r(i)\leq e_{j+1}$. Since there is no element of $E$ strictly between them, $r(i)$ is one of these two endpoints. Moreover the indices mapped left must all precede the indices mapped right. There is therefore a single cut in the gap. Sign preservation is equivalent to the signs before the cut being $\varepsilon_{e_j}$ and those after it being $\varepsilon_{e_{j+1}}$. This proves necessity of the gap condition. If the endpoint signs differ, an admissible cut is unique. If they agree, every sign in the gap must have that common value and all $d+1$ cuts are admissible, where $d=e_{j+1}-e_j-1$.

Conversely, choose an admissible cut independently in every gap, use the forced exterior values, and fix every $e_j$. The resulting map is nondecreasing across each gap and its boundary, hence globally nondecreasing; it preserves signs by construction and retracts onto $E$. These choices describe every retraction exactly once. Multiplying their independent numbers gives the displayed product. The signed-edge correspondence turns each retraction back into a unique section of $p_E$, completing the proof.

For $\tau=+-$ and $E=\{1\}$ the trailing negative sign violates the exterior condition. Thus the surjective vertex map $+ -\to+$ has no section. For $\tau=+++$ and $E=\{1,3\}$ the gap has one omitted positive edge, and there are exactly two sections; their vertex maps $[0,2]\to[0,3]$ are $(0,1,3)$ and $(0,2,3)$. This distinguishes the collapse of an edge from a simplicial degeneracy with an available section.

We checked the ordinal description and comonoid interpretation in \citet[Appendix~D]{AguiarMahajan2010}, the zig-zag and Reedy structures of \citet[Definition~1.1 and Example~4.12]{ArakawaCnossen2026}, and searched for signed or coloured ordinal retractions and sections of Grothendieck's interval maps. We could not find this exact interval-form criterion and count in those sources. The result is an elementary consequence of the normal form, recorded here as a proved continuation of the calculation, without a claim of priority.
\end{dcomment}

\begin{dcomment}{context}{p49}{Concatenation is a pushout with chosen endpoints}
Grothendieck's reason for introducing types now becomes visible: the concatenation of words should encode the actual gluing of paths, including their transformations. An interval is a free category on an oriented linear graph. Gluing the last vertex of the first graph to the first vertex of the second produces the graph of the concatenated word, and the free-category construction preserves this pushout. Therefore
\[
 I_{\tau\tau'}=I_\tau\amalg_e I_{\tau'}
\]
is a pushout in $\Cat$ with the specified endpoint maps. Any directed path between two vertices of one summand stays in that summand: entering the other would require crossing the gluing vertex twice in an underlying linear graph. Thus the inclusions are fully faithful as well as injective on objects.

The corresponding assertion about path categories is also literal. Restriction gives an isomorphism
\[
 \operatorname{Fun}(I_{\tau\tau'},X)
 \cong
 \operatorname{Fun}(I_\tau,X)
 \times_X\operatorname{Fun}(I_{\tau'},X),
\]
where the fibre product uses evaluation at the two endpoints being identified. On objects, two functors agree at the common vertex exactly when their two paths concatenate. On morphisms, two natural transformations glue exactly when their components at that vertex are equal. Naturality along each generating edge then gives naturality on every composite. This verifies the categorical pullback, rather than only the corresponding assertion about sets of paths.

The convention on the order of composition now has a concrete meaning. The word $\tau\tau'$ is read first through $\tau$ and then through $\tau'$, whereas the corresponding categorical composition is written $c'\circ c$. This harmless reversal of written order should be fixed before discussing fractions. It is not an opposite-category operation and does not reverse any transition.

The endpoint data are indispensable. The parenthetical assertion on manuscript page~50 that concatenation also gives a bifunctor on $\uTyp$, as defined on page~30 using \emph{all} functors, is false. The one-arrow categories $I_+$ and $I_-$ are isomorphic in $\uTyp$. A bifunctor whose object operation is word concatenation would carry this isomorphism, together with $\mathrm{id}_{I_+}$, to an isomorphism between $I_{++}$ and $I_{-+}$. The former is a three-object chain; the latter has one initial object and two incomparable maximal objects. They are not isomorphic. Thus the claimed bifunctor cannot exist on that larger category. It does exist on endpoint-preserving functors, even without monotonicity, because their restrictions agree at the glued vertex. The final category $\uTy$ satisfies this requirement and remains closed under concatenation. This correction isolates precisely the structure that the path formalism needs.
\end{dcomment}

\begin{dcomment}{later}{p50}{The type category is generated by two comonoids}
The strict monoidal structure places Grothendieck's calculation beside a classical universal property of the simplicial category. In the single-sign case, endpoint-preserving maps of finite intervals are opposite to monotone maps of finite ordered edge sets. The latter category, with ordinal sum, contains the universal monoid; its opposite contains the universal counital coassociative comonoid. An explicit account, with the weak-composition coordinates used here, is \citet[Appendix~D.1--D.2, especially Proposition~D.2, pp.~713--716]{AguiarMahajan2010}, which traces the universal property to Mac Lane's Proposition~VII.5.1. The two signs simply keep two independent copies of this structure in the same ordered word.

In $\uTy$ put
\[
 \delta_\varepsilon=\mathfrak f_{\varepsilon,2}:
      \Delta_\varepsilon\longrightarrow\Delta_\varepsilon^2,
 \qquad
 e_\varepsilon=\mathfrak f_{\varepsilon,0}:
      \Delta_\varepsilon\longrightarrow1,
 \qquad \varepsilon\in\{+,-\}.
\]
Both ways of applying $\delta_\varepsilon$ twice give the unique endpoint-preserving map from $\Delta_\varepsilon$ to $\Delta_\varepsilon^3$. Applying $e_\varepsilon$ to either output of $\delta_\varepsilon$ gives the identity, since its surviving weight is $1$. Thus each $\Delta_\varepsilon$ is a counital coassociative comonoid. There is no specified relation between the positive and negative comonoids.

More precisely, a strict monoidal functor from $\uTy$ to a strict monoidal category $(\mathcal C,\otimes,\mathbf1)$ is uniquely specified by a pair of counital coassociative comonoids $C_+$ and $C_-$ in $\mathcal C$. Here is the construction, which also proves the assertion. For either comonoid, let $\delta^{(a)}:C\to C^{\otimes a}$ denote the iterated comultiplication, with $\delta^{(0)}$ its counit and $\delta^{(1)}=\mathrm{id}$. Coassociativity and the counit laws give, for every list $b_1,\ldots,b_a$ of nonnegative integers,
\[
 \bigl(\delta^{(b_1)}\otimes\cdots\otimes\delta^{(b_a)}\bigr)
       \delta^{(a)}
 =\delta^{(b_1+\cdots+b_a)}.
\]
For completeness, the binary identity
$(\delta^{(p)}\otimes\delta^{(q)})\delta^{(2)}
 =\delta^{(p+q)}$ follows by induction on positive $p,q$ using coassociativity; the cases $p=0$ or $q=0$ are the counit identities, including $p=q=0$. Applying this binary identity successively to the first $a-1$ factors and the last factor proves the displayed formula by induction on $a$. The case $a=0$ says that the identity of the unit composed with the counit is the counit.

Send the word $\varepsilon_1\cdots\varepsilon_n$ to
$C_{\varepsilon_1}\otimes\cdots\otimes C_{\varepsilon_n}$ and send its outgoing morphism of weights $(a_1,\ldots,a_n)$ to
\[
 \delta_{\varepsilon_1}^{(a_1)}\otimes\cdots\otimes
 \delta_{\varepsilon_n}^{(a_n)}.
\]
The block-sum formula for composition and the preceding identity prove that this assignment is a functor. The all-one vector gives identities; the empty word gives the unit; concatenation of vectors gives the tensor product of maps. It is therefore strict monoidal. Conversely every outgoing morphism is the tensor product of its one-edge factors, and each such factor is an iterated comultiplication or counit. A strict monoidal functor must have the displayed values, establishing uniqueness.

This is a reinterpretation of the known ordinal/comonoid universal property, not a new homotopy theorem. It explains the two apparently different generators in the manuscript: $\mathfrak f_{\varepsilon,0}$ erases a transition, whereas $\mathfrak f_{\varepsilon,2}$ subdivides one. Their iterates produce every allowed reparametrization. No braiding is present: the order of the source transitions is part of the structure and cannot be permuted.
\end{dcomment}

\begin{dcomment}{exposition}{p53}{Free monoids in degrees one and two}
Grothendieck uses the nerve to keep composition together with concatenation. A string of composable reparametrizations can still be cut along the original source edges. The resulting factors are independent, and this is the reason the monoid in each nerve degree is free.

In degree one, a morphism with source word $\varepsilon_1\cdots\varepsilon_n$ and weights $(a_1,\ldots,a_n)$ has the unique decomposition
\[
 f=\mathfrak f_{\varepsilon_1,a_1}*\cdots*
       \mathfrak f_{\varepsilon_n,a_n}.
\]
The factors are indexed by source edges, even when their weights vanish. A generator $\mathfrak f_{\varepsilon,0}$ is therefore not the unit of this free monoid: its source has one edge. The monoid unit is the empty factorization, namely the identity of the point. Nor is $\mathfrak f_{\varepsilon,1}$ the monoid unit; it is the identity of a one-edge object and is indecomposable under concatenation. This separates the categorical identity operation from the unit of the tensor product.

Now consider $\tau\xrightarrow f\tau'\xrightarrow g\tau''$. The $i$th original edge produces an ordered block of $a_i$ edges after $f$. Let $\pi_i:[1,a_i]\to\mathbb N$ record the weights that $g$ assigns to that block. All descendants have the original sign $\varepsilon_i$. Consequently the pair has the unique factorization
\[
 (f,g)=\prod_{i=1}^{n,*}\mathfrak f_{\varepsilon_i,a_i,\pi_i}.
\]
Existence comes from cutting every intermediate word into its descendant blocks; uniqueness comes from the fixed order of the original edges and the unique weight vector of each restriction. The nullary case $a_i=0$ has one empty function $\pi_i$, so it contributes a legitimate generator whose two targets are the point. For $n=0$ the whole pair is the identity string at the point and is the monoid unit. This proves the claimed freeness in degree two without discarding extinct blocks.

Writing $N_r=N_r(\uTy)$, the three faces of a degree-two generator are
\[
\begin{aligned}
 d_2(\mathfrak f_{\varepsilon,a,\pi})
    &=\mathfrak f_{\varepsilon,a},\\
 d_1(\mathfrak f_{\varepsilon,a,\pi})
    &=\mathfrak f_{\varepsilon,\sum_{j=1}^{a}\pi(j)},\\
 d_0(\mathfrak f_{\varepsilon,a,\pi})
    &=\mathfrak f_{\varepsilon,\pi(1)}*\cdots*
        \mathfrak f_{\varepsilon,\pi(a)}.
\end{aligned}
\]
The last line is the empty product when $a=0$. The two degeneracies of a degree-one generator insert the identity before or after its map:
\[
 s_0(\mathfrak f_{\varepsilon,a})
       =\mathfrak f_{\varepsilon,1,(a)},\qquad
 s_1(\mathfrak f_{\varepsilon,a})
       =\mathfrak f_{\varepsilon,a,1_a}.
\]
Here $(a)$ is the one-entry weight vector and $1_a$ the all-one vector of length $a$. These formulas recover precisely the composition and degeneracies described on manuscript pages~54--56, while displaying the other two faces as well.

The word ``semi-simplicial'' in this edition includes degeneracies; in current terminology the nerve is a simplicial monoid. The faces and degeneracies are monoid homomorphisms because concatenation is a bifunctor and is strictly associative and unital. This statement concerns the whole degreewise free monoid. Its distinguished basis need not itself be a simplicial set: $d_0$ can send one generator to a product of several generators, or to the unit.
\end{dcomment}

\begin{dcomment}{exposition}{p57}{All nerve degrees as ordered forests}
The final assertion of manuscript page~57 is that the same decomposition works for arbitrarily long strings of reparametrizations. It is useful to complete the argument now, since it explains both the support conditions at the bottom of the page and the one face map that behaves differently from the others.

Fix $r\geq0$. Call an $r$-step string \emph{rooted} when its initial type has one transition. Every intermediate type is then a power of the same sign, and the string has the form
\[
 \Delta_\varepsilon^{l_0}\longrightarrow
 \Delta_\varepsilon^{l_1}\longrightarrow\cdots\longrightarrow
 \Delta_\varepsilon^{l_r},\qquad l_0=1.
\]
For $0\leq a<r$ its $a$th weight vector is
\[
 \pi_a:[1,l_a]\longrightarrow\mathbb N,
 \qquad l_{a+1}=\sum_{j=1}^{l_a}\pi_a(j).
\]
One can regard this as an ordered rooted tree with $r$ levels of edges: a transition at level $a$ has $\pi_a(j)$ ordered descendants at level $a+1$. A vertex may have no descendants. Once a whole level is empty, all subsequent levels are empty because there is no morphism from the point to a nonpoint type.

For a general string, let its initial word have transitions $1,\ldots,n$. At the first stage the descendants of each original transition form a consecutive monochromatic block. At the next stage the descendants of all transitions in one such block again form a consecutive block: concatenate their individual image blocks, allowing zero lengths. Induction gives such a block at every level. Restricting all maps to the descendants of a fixed original edge produces a rooted string, and concatenating the $n$ rooted strings reconstructs the original string. Each boundary is determined by cumulative weights, so this factorization is unique. A rooted string cannot be a product of two nonunit strings, since the lengths of their initial words would be positive integers summing to one. Conversely a string whose initial word has at least two edges splits after its first edge. Thus rooted strings are exactly the indecomposables, and
\[
 N_r(\uTy)\cong\text{the free monoid on rooted $r$-step strings}.
\]
For $r=0$ there are exactly two generators, the two signs. For an empty initial word there is only the all-identity string at the point, which is the monoid unit. These boundary cases establish the assertion for every $r$.

The simplicial operators can now be read directly on these forests. An inner face removes one level and replaces the two successive weight vectors by the sums over consecutive descendant blocks. The last face forgets the last level. A degeneracy repeats a level, inserting weight $1$ at each of its transitions. These operations preserve a single root. The first face removes the original root level: a rooted string with $l_1$ first-generation transitions becomes the concatenation of their $l_1$ rooted descendant strings. This is the higher-dimensional version of the formula for $d_0$ above. Since these are exactly restriction and composition of actual functors $[r]\to\uTy$, they satisfy all simplicial identities; the unique forest decomposition identifies their formulas with those categorical operators.

For $r>0$, to use fixed-domain functions as on manuscript page~57, extend each $\pi_a$ by zero to $\mathbb N^*$. The exact conditions are
\[
 \operatorname{supp}(\pi_0)\subseteq\{1\},\qquad
 \operatorname{supp}(\pi_a)\subseteq
       [1,\operatorname{long}(\pi_{a-1})]
       \quad(1\leq a<r).
\]
The source's final boxed range starts at $a=0$, although $\pi_{-1}$ is undefined; it must begin at $1$. The preceding display defining $\operatorname{long}(\pi_{i-1})$ likewise needs the indices $j\in[1,l_{i-1}]$ and $\pi_{i-1}(j)$. These printed indexing slips do not affect the argument once the successive domains are kept in view.

One should not infer from this result that the simplicial monoid is freely generated by a simplicial set of indecomposables: the first-face calculation rules that out for this particular basis. Nor does the algebraic freeness make its classifying space complicated. The point is terminal in $\uTy$, since every weight vector admits the unique further collapse to zero. Appending that terminal object gives the usual contraction of the nerve. The free presentations describe how all reparametrizations fit together, while this terminal object describes their ordinary homotopy type.
\end{dcomment}

\begin{dcomment}{later}{p31:Je vais exiger que les morphismes}{The same type category in later zigzag constructions}
Grothendieck's final restriction on changes of shape has a precise counterpart in the later literature on localization. The correspondence can be checked directly, before making any assertion about homotopy theory. Given a word $\tau$, mark the backward transitions of $I_\tau$ and their composites, together with the identities. This makes $I_\tau$ a relative category. An increasing endpoint-preserving functor automatically respects this marking: a nonconstant image of a backward transition is a composite of backward transitions. Conversely, forgetting the marking of an increasing endpoint-preserving relative functor gives exactly a morphism of $\uTy$. The objects, morphisms, and concatenation are therefore unchanged by this decoration.

Under $+\leftrightarrow A$ and $-\leftrightarrow W^{-1}$, this identifies $\uTy$ with the category $Z$ of \citet[Definition~3.5, Remark~3.6, and footnotes~8--10, pp.~341--342]{MazelGee2018Hammocks}. Their weak-equivalence labels restrict which diagrams in a coefficient category are allowed; they do not impose an additional restriction on these monotone shape maps. The empty word remains the monoidal unit. This is an identification of the finite indexing categories, rather than an inference from a resemblance between pictures.

The image factorization implicit in the manuscript also explains a useful structure. If $f:I_\tau\to I_\upsilon$ is increasing, the full subcategory on its image is an interval. Factor $f$ through its standard oriented model:
\[
 I_\tau\mathrel{\mathop{\longrightarrow}^{q}}I_\omega
 \mathrel{\mathop{\lhook\joinrel\longrightarrow}^{j}}I_\upsilon.
\]
The first map is surjective on vertices and the second injective. Their uniqueness follows from the unique increasing enumeration of the image. Nonidentity surjections lower the length, and nonidentity injections raise it. Thus length is a Reedy degree. The reduced interval gives the smallest image of an endpoint-preserving injection: Proposition~5 says that every such image contains all critical vertices. If the interval is not reduced, its reduced inclusion is consequently initial among the proper injections into it; for a reduced interval there are none. These statements follow from the full-subcategory argument already given, without a new hypothesis on coefficients.

This Reedy structure and this exact use of critical vertices occur in \citet[Example~4.12 and Lemma~4.16]{ArakawaCnossen2026}, where the category is denoted $\mathrm{Zig}_{\partial}$. Their work is an August~2026 preprint. We cite these known applications, rather than present the Reedy interpretation as an original result. What remains to be added to the present manuscript is the relative and homotopical diagram theory over this category of shapes.
\end{dcomment}

\begin{dcomment}{status}{p57}{What the shape calculus does for localization}
The free monoids just constructed organize subdivision and concatenation. To obtain mapping spaces, a second ingredient is needed: comparisons of labelled paths must carry the specified weak-equivalence information. For a relative category $(C,W)$ and fixed $x,y$, let $Z_\tau(C,W)_{x,y}$ be the category of paths of type $\tau$ from $x$ to $y$, with backward arrows in $W$, and natural transformations whose components lie in $W$ and whose endpoint components are identities. Precomposition makes these categories contravariant in $\uTy$. Their nerves, or their groupoid completions, retain comparisons that the set of path components alone would discard.

The analogous construction for relative infinity-categories is developed by \citet[Definitions~3.9 and~3.17]{MazelGee2018Hammocks}. Under a homotopical three-arrow calculus, his Theorem~4.4 identifies the geometric realizations of its hammock mapping objects with the groupoid completions of three-arrow zigzag categories. This is an additional hypothesis, not a consequence of the finite normal form in the manuscript. Independently, \citet[Theorem~3.8]{MazelGee2019Rezk} proves that completing the relative Rezk nerve presents localization. That universal property supplies the criterion with which a proposed path model must be compared.

For the general mapping-space comparison, \citet[Theorem~1.9 and Corollary~1.11]{ArakawaCnossen2026} prove in their preprint that
\[
 \operatorname{Map}_{C[W^{-1}]}(x,y)
 \simeq \operatorname*{colim}_{\tau\in\uTy^{\mathrm{op}}}
          \bigl|Z_\tau(C,W)_{x,y}\bigr|,
\]
where the colimit is taken in spaces and $|{-}|$ is groupoid completion. The displayed formula uses the ordinary relative categories just defined; the passage from their strict endpoint fibres to the invariant formula is supplied by Convention~4.2, Definition~4.3, Theorem~4.7, and Proposition~4.8(b) of the same paper. The general infinity-categorical statement uses the corresponding infinity-categories of relative diagrams and their endpoint fibres. Their Corollary~5.7 supplies the localization identification for Mazel-Gee's hammock construction. These references distinguish the established Rezk-nerve theorem from the more recent explicit comparison.

The comparison makes the scope of the present pages clear. They provide the shapes, their functorial changes, and strict concatenation. They do not yet choose $W$ or impose the conditions on vertical arrows used above. No new open conjecture is required to finish their finite combinatorics. The subsequent localization construction still has to be checked with its own morphisms and universal property; the contractibility of an interval, by itself, does not perform that check.
\end{dcomment}
